\section{Introdução}\index{espaço euclidiano}
Um ponto $p \in \mathbb R^n$ é descrito por $n$ coordenadas reais.
\[
    p = (x_1, \dots, x_n).
\]

Lembremos que $\mathbb R^2$ é o conjunto de todos os pares ordenados $(x, y)$, onde $x, y \in \mathbb R$.
Em símbolos:
\begin{equation*}
    \mathbb R^2 = \{(x, y) : x, y \in \mathbb R\}.
\end{equation*}

Em analogia à representação de $\mathbb R$ como uma reta, podemos visualizar $\mathbb R^2$ como o conjunto dos pontos de um plano cartesiano.

Na Figura~\ref{fig:plano-cartesiano}, o ponto $(2,3)$ está posicionado $2$ unidades no sentido positivo do eixo $x$ e $3$ unidades no sentido positivo do eixo $y$.

\begin{figure}[ht]
    \centering
    \begin{tikzpicture}
    \draw[->,ultra thick] (-4,0)--(4,0) node[right]{$x$};
    \draw[->,ultra thick] (0,-4)--(0,4) node[above]{$y$};
    % Ponto (2,3)
    \filldraw[black] (2,3) circle (2pt);
    \node[above right] at (2,3) {$(2,3)$};
    \draw[dotted] (2,0) -- (2,3);
    \draw[dotted] (0,3) -- (2,3);
    \draw (2,0.1) -- (2,-0.1);
    \node[below] at (2,0) {2};
    \draw (0.1,3) -- (-0.1,3);
    \node[left] at (0,3) {3};
    \node[below left] at (0,0) {0};

    \end{tikzpicture}
    \caption{O plano cartesiano.}
    \label{fig:plano-cartesiano}
\end{figure}

Por sua vez, o conjunto de todas as triplas ordenadas de números reais é denotado por $\mathbb R^3$.
Em símbolos:
\begin{equation*}
    \mathbb R^3 = \{(x, y, z) : x, y, z \in \mathbb R\}.
\end{equation*}

Seguindo o padrão já comentado, podemos visualizar $\mathbb R^3$ como o conjunto dos pontos do espaço tridimensional.

A tripla $(x,y,z)$ representa um ponto com três coordenadas: a primeira associada ao eixo $x$, a segunda ao eixo $y$ e a terceira ao eixo $z$.
O ponto $(2,3,1)$ está representado na Figura~\ref{fig:espaco-tridimensional}.

\begin{figure}[ht]
    \centering
    \tdplotsetmaincoords{70}{110}
    \begin{tikzpicture}[tdplot_main_coords]
    % Eixos
    \draw[->,ultra thick] (-4,0,0)--(4,0,0) node[right]{$x$};
    \draw[->,ultra thick] (0,-4,0)--(0,4,0) node[above]{$y$};
    \draw[->,ultra thick] (0,0,-4)--(0,0,4) node[above]{$z$};
    % Ponto (2,3,1)
    \filldraw[black] (2,3,1) circle (2pt);
    \node[above right] at (2,3,1) {$(2,3,1)$};
    % Linhas pontilhadas
    \draw[dotted] (2,0,0) -- (2,3,0) -- (2,3,1);
    \draw[dotted] (0,3,0) -- (2,3,0);
    \draw[dotted] (2,0,0) -- (2,0,1) -- (2,3,1);
    \draw[dotted] (0,0,1) -- (2,0,1);
    % Marcas nos eixos
    \draw (2,0.1,0) -- (2,-0.1,0);
    \node[below] at (2,0,0) {2};
    \draw (0.1,3,0) -- (-0.1,3,0);
    \node[left] at (0,3,0) {3};
    \draw (0.1,0,1) -- (-0.1,0,1);
    \node[left] at (0,0,1) {1};
    \node[below left] at (0,0,0) {0};

    \end{tikzpicture}
    \caption{O espaço tridimensional.}
    \label{fig:espaco-tridimensional}
\end{figure}

De modo geral, para cada inteiro $n\geq 1$, o conjunto $\mathbb R^n$ é o conjunto de todas as $n$-tuplas ordenadas de números reais:
\begin{equation*}
    \mathbb R^n = \{(x_1, \ldots, x_n) : x_1, \ldots, x_n \in \mathbb R\}.
\end{equation*}

Para $n\geq 4$, esse conjunto não possui uma representação gráfica simples no estilo dos casos anteriores, de dimensões menores.
Porém, a teoria desenvolvida para $\mathbb R^n$, com $n\geq 4$, é uma extensão natural da teoria desenvolvida para $\mathbb R$, $\mathbb R^2$ e $\mathbb R^3$, e possui amplas aplicações práticas e teóricas.

Cada espaço $\mathbb R^n$ possui um elemento denominado \emph{origem}\index{origem}, denotado por $\vec 0$, ou simplesmente por $0$, cujas coordenadas são todas iguais ao número real $0$.
\begin{equation*}
    \vec 0 = (0, \ldots, 0).
\end{equation*}

Um detalhe que pode ser confuso à primeira vista, mas que não gera confusão, na prática, é que a notação $0$ é utilizada tanto para o número real $0$ quanto para a origem $\vec 0$ de $\mathbb R^n$.
Ao se escrever $0=(0, \ldots, 0)$, o primeiro símbolo $0$ representa a origem de $\mathbb R^n$, enquanto os demais representam o número real $0$, sendo, portanto, objetos de natureza diferente.

No caso $n=1$, identificamos usualmente a $1$-tupla $(x)$ com o próprio número real $x$, de modo que $\mathbb R^1$ é identificado com $\mathbb R$.
Com essa identificação, a origem de $\mathbb R^1$ é o número real $0$.

Elementos de $\mathbb R^n$ serão frequentemente chamados de \emph{pontos}\index{ponto} ou \emph{vetores}\index{vetor}.
Essa dupla nomenclatura reflete duas formas diferentes de interpretar o mesmo objeto matemático.
A palavra \emph{ponto} será usualmente utilizada em situações em que se faz referência a posições no espaço, enquanto a palavra \emph{vetor} é mais utilizada em situações em que pensamos no segmento orientado que parte da origem $\vec 0=(0, \ldots, 0)$\index{origem} e termina no ponto em questão.
Quando esse ponto não é a origem, esse segmento determina, intuitivamente, uma direção, um sentido e um comprimento.
Quando esse ponto é a origem, obtemos o \emph{vetor nulo}\index{vetor nulo}, que possui comprimento zero e não determina uma direção ou um sentido.

O conjunto $\mathbb R^n$, munido das operações e da distância usuais que estudaremos a seguir, é chamado de \emph{espaço euclidiano} $n$-dimensional.